\documentclass[11pt,a4paper]{scrartcl}

\usepackage{../manostilius_lt}

\title{Daugianarių dalyba}
\author{Andrius Berniukevičius}

\begin{document}

\maketitle

\section{Daugianario dalyba iš dvinario}

\begin{definition}
Tegu $P(x)$ ir $D(x)$ yra daugianariai, o $D(x) \ne 0$. Jei egzistuoja daugianaris $Q(x)$, kuriam
\[
P(x)=D(x)\cdot Q(x),
\]
Daugianaris $P(x)$ \vocab{dalijasi} iš daugianario $D(x)$, o daugianaris $Q(x)$ vadinamas \vocab{dalmeniu}.
\end{definition}

\begin{algorithm}
Norėdami padalyti daugianarį iš dvinario stulpeliu, atliekame šiuos veiksmus:
\begin{enumerate}
    \item Padalijame dalinio vyriausiąjį narį iš daliklio vyriausiojo nario. Gautą narį rašome į dalmenį.
    \item Naują dalmens narį padauginame iš viso daliklio ir gautą daugianarį rašome po daliniu, sutapatindami vienodų laipsnių narius.
    \item Atimame gautą sandaugą iš dalinio.
	\item Kartojame veiksmus, kol gauname nulį.
\end{enumerate}
\end{algorithm}

\section{Pavyzdžiai}

\begin{example}
Padalykime $6x^4-x^3-22x^2+17x-3$ iš $2x-3$.

\setlength{\arraycolsep}{2pt}
\renewcommand{\arraystretch}{1.1}
\[
\begin{array}[t]{ccccccccc}
  6x^4 & - & x^3  & - & 22x^2 & + & 17x & - & 3 \\
  6x^4 & - & 9x^3 &   & \downarrow & & & & \\
  \cline{1-3}
  \noalign{\vskip 3pt}
	  &   & 8x^3  & - & 22x^2 &   & \downarrow & & \\
	  &   & 8x^3  & - & 12x^2 &   &             & & \\
  \cline{3-5}
  \noalign{\vskip 3pt}
	  &   &       &   & -10x^2 & + & 17x &   & \downarrow \\
	  &   &       &   & -10x^2 & + & 15x &   & \\
  \cline{5-7}
  \noalign{\vskip 3pt}
	  &   &       &   &        &   & 2x & - & 3 \\
	  &   &       &   &        &   & 2x & - & 3 \\
  \cline{7-9}
  \noalign{\vskip 3pt}
	  &   &       &   &        &   &    &   & 0
\end{array}
\enskip
\begin{array}[t]{|@{\enskip}l}
  2x-3 \\
  \hline
  \rule{0pt}{14pt}3x^3+4x^2-5x+1
\end{array}
\]
Dalmenį sudarome po vieną narį:
\begin{enumerate}
	\item Dalinio vyriausiąjį narį dalijame iš daliklio vyriausiojo nario:
	\[
	6x^4:(2x)=3x^3.
	\]
	\item Padauginame $3x^3(2x-3)=6x^4-9x^3$ ir atimame. Gauname $8x^3$, tuomet nuleidžiame $-22x^2$, todėl gauname $8x^3-22x^2$. Dabar
	\[
	8x^3:(2x)=4x^2.
	\]
	\item Padauginame $4x^2(2x-3)=8x^3-12x^2$ ir atimame. Gauname $-10x^2$, tuomet nuleidžiame $+17x$, todėl gauname $-10x^2+17x$. Dabar
	\[
	-10x^2:(2x)=-5x.
	\]
	\item Padauginame $-5x(2x-3)=-10x^2+15x$ ir atimame. Gauname $2x$, tuomet nuleidžiame $-3$, todėl gauname $2x-3$. Dabar
	\[
	2x:(2x)=1.
	\]
\end{enumerate}
Taigi
\[
(6x^4-x^3-22x^2+17x-3):(2x-3)=3x^3+4x^2-5x+1.
\]
\end{example}

\begin{example}
Padalykime $2x^3+x^2-11x+2$ iš $x-2$.

\setlength{\arraycolsep}{2pt}
\renewcommand{\arraystretch}{1.1}
\[
\begin{array}[t]{ccccccc}
  2x^3 & + & x^2  & - & 11x & + & 2 \\
  2x^3 & - & 4x^2 &   & \downarrow & & \\
  \cline{1-3}
  \noalign{\vskip 3pt}
	  &   & 5x^2  & - & 11x &   & \downarrow \\
	  &   & 5x^2  & - & 10x &   & \\
  \cline{3-5}
  \noalign{\vskip 3pt}
	  &   &       &   & -x  & + & 2 \\
	  &   &       &   & -x  & + & 2 \\
  \cline{5-7}
  \noalign{\vskip 3pt}
	  &   &       &   &     &   & 0
\end{array}
\enskip
\begin{array}[t]{|@{\enskip}l}
  x-2 \\
  \hline
  \rule{0pt}{14pt}2x^2+5x-1
\end{array}
\]
Dalmenį sudarome po vieną narį:
\begin{enumerate}
	\item Dalinio vyriausiąjį narį dalijame iš daliklio vyriausiojo nario:
	\[
	2x^3:x=2x^2.
	\]
	\item Padauginame $2x^2(x-2)=2x^3-4x^2$ ir atimame. Gauname $5x^2$, tuomet nuleidžiame $-11x$, todėl gauname $5x^2-11x$. Dabar
	\[
	5x^2:x=5x.
	\]
	\item Padauginame $5x(x-2)=5x^2-10x$ ir atimame. Gauname $-x$, tuomet nuleidžiame $+2$, todėl gauname $-x+2$. Dabar
	\[
	-x:x=-1.
	\]
\end{enumerate}
Vadinasi,
\[
(2x^3+x^2-11x+2):(x-2)=2x^2+5x-1.
\]
\end{example}

\begin{example}
Padalykime $x^3+2x^2+4x+21$ iš $x+3$.

\setlength{\arraycolsep}{2pt}
\renewcommand{\arraystretch}{1.1}
\[
\begin{array}[t]{ccccccc}
	x^3 & + & 2x^2 & + & 4x & + & 21 \\
	x^3 & + & 3x^2 &      & \downarrow &   &    \\
	\cline{1-3}
	\noalign{\vskip 3pt}
	    &   & -x^2 & + & 4x &   & \downarrow \\
	    &   & -x^2 & - & 3x &   &             \\
	\cline{3-5}
	\noalign{\vskip 3pt}
	    &   &      &   & 7x & + & 21 \\
	    &   &      &   & 7x & + & 21 \\
	\cline{5-7}
	\noalign{\vskip 3pt}
	    &   &      &   &    &   & 0
\end{array}
\enskip
\begin{array}[t]{|@{\enskip}l}
	x+3 \\
	\hline
	\rule{0pt}{14pt}x^2-x+7
\end{array}
\]
Dalmenį sudarome po vieną narį:
\begin{enumerate}
	\item Dalinio vyriausiąjį narį dalijame iš daliklio vyriausiojo nario:
	\[
	x^3:x=x^2.
	\]
	\item Padauginame $x^2(x+3)=x^3+3x^2$ ir atimame. Gauname $-x^2$, tuomet nuleidžiame $+4x$, todėl gauname $-x^2+4x$. Dabar
	\[
	-x^2:x=-x.
	\]
	\item Padauginame $-x(x+3)=-x^2-3x$ ir atimame. Gauname $7x$, tuomet nuleidžiame $+21$, todėl gauname $7x+21$. Dabar
	\[
	7x:x=7.
	\]
\end{enumerate}
Taigi
\[
(x^3+2x^2+4x+21):(x+3)=x^2-x+7.
\]
\end{example}

\section{Uždaviniai}

\begin{problem}
Atlikite daugianarių dalybą ir nurodykite dalmenį.
\begin{tasks}[label={\arabic*.},label-offset=0.9em](2)
	\task $(x^3+5x^2+8x+4):(x+1)$
	\task $(2x^3-3x^2-8x+12):(x-2)$
	\task $(3x^3+x^2-10x-24):(x+2)$
	\task $(2x^4-5x^3-3x^2+20x-20):(x-2)$
	\task $(x^3-4x+3):(x-1)$
	\task $(4x^3-7x^2-7x+10):(x-2)$
	\task $(x^4+4x^3+x^2-4x+4):(x+2)$
	\task $(5x^4-3x^3-3x^2+5x-4):(x-1)$
	\task $(x^4+3x^3+x^2-3x-2):(x^2-1)$
	\task $(2x^4+5x^3-7x^2+10x-4):(2x-1)$
\end{tasks}
\end{problem}

\newpage
\section{Atsakymai}

\begin{tasks}[label={\arabic*.},label-offset=0.9em](1)
	\task $x^2+4x+4$
	\task $2x^2+x-6$
	\task $3x^2-5x$
	\task $2x^3-x^2-5x+10$
	\task $x^2+x-3$
	\task $4x^2+x-5$
	\task $x^3+2x^2-3x+2$
	\task $5x^3+2x^2-x+4$
	\task $x^2+3x+2$
	\task $x^3+3x^2-2x+4$
\end{tasks}

\end{document}
